% !TeX program = pdflatex
% Compiler Engineering Suite - shared standalone pdfLaTeX preamble
\documentclass[10pt,letterpaper]{article}
\usepackage[T1]{fontenc}
\usepackage[utf8]{inputenc}
\usepackage{lmodern}
\usepackage[margin=.78in,headheight=18pt]{geometry}
\usepackage[table]{xcolor}
\usepackage{microtype}
\usepackage{amsmath,amssymb}
\usepackage{booktabs,longtable,array}
\usepackage{enumitem}
\usepackage{fancyhdr}
\usepackage{hyperref}
\usepackage{needspace}
\usepackage{fancyvrb}
\usepackage{tikz}
\usetikzlibrary{arrows.meta,positioning,shapes.geometric}
\definecolor{CompilerBlue}{HTML}{1F4E79}
\definecolor{CompilerInk}{HTML}{263746}
\definecolor{CompilerGray}{HTML}{52616B}
\definecolor{CompilerLight}{HTML}{EFF4F8}
\hypersetup{colorlinks=true,linkcolor=CompilerBlue,urlcolor=CompilerBlue,citecolor=CompilerBlue}
\setlength{\parindent}{0pt}
\setlength{\parskip}{6pt plus 2pt}
\setlist[itemize]{leftmargin=*,itemsep=3pt,topsep=2pt}
\setlist[enumerate]{leftmargin=*,itemsep=3pt,topsep=2pt}
\pagestyle{fancy}
\fancyhf{}
\fancyhead[L]{\small\textcolor{CompilerGray}{Compiler Engineering Documentation Suite}}
\fancyhead[R]{\small\thepage}
\renewcommand{\headrulewidth}{.4pt}
\newcommand{\evidence}[1]{\par\smallskip\noindent\textbf{Evidence required:} #1\par}
\newcommand{\sourceanchor}[1]{\textbf{Book anchor:} #1\par}
\newcommand{\modern}{\textbf{Modern engineering extension (not attributed to the source book):} }
\newcommand{\writinglines}[1]{\par\noindent\rule{\linewidth}{.3pt}\par\vspace{#1\baselineskip}\noindent\rule{\linewidth}{.3pt}\par}
\newcommand{\checkboxline}[1]{\noindent $\square$\quad #1\par}
\newcommand{\formfield}[2]{\noindent\textbf{#1}\par\vspace{#2\baselineskip}\hrule\medskip}

\hypersetup{pdftitle={Chapter 8: Introduction to Optimization},pdfauthor={Compiler Engineering Documentation Suite}}
\begin{document}
\begin{titlepage}
\centering\vspace*{1.0in}
{\Huge\bfseries\textcolor{CompilerBlue}{Chapter 8: Introduction to Optimization}\par}
\vspace{.22in}{\Large Critical reading, original exercises, evidence and review\par}
\vspace{.4in}{\large Compiler Engineering Documentation Suite\par}
\vfill
\begin{minipage}{.86\linewidth}\small
Source basis: Keith D. Cooper and Linda Torczon, \emph{Engineering a Compiler}, Third Edition (Morgan Kaufmann/Elsevier, 2023). All explanatory text, examples, questions, and implementation designs in this document are independently authored. The book is not included.\par\medskip
Document class: Chapter study guide.
\end{minipage}
\vfill\textcolor{CompilerGray}{Version 1.0 --- 2026-10-30}\par
\end{titlepage}
\tableofcontents\clearpage
\section{Scope and source crosswalk}
\sourceanchor{Chapter 8: Introduction to Optimization; sections 8.2--8.7; pp. 379--444.}
\textbf{Core topics:} Local, regional, global and interprocedural optimization.

This guide records the chapter learning sequence as an independent engineering practice plan; consult the source text for proofs, figures, terminology and its own exercises.

\section{Competency outcomes}
\begin{itemize}
\item Build a legal optimization pass and distinguish soundness from profitability.
\item Define value numbering, constant propagation, cloning, inlining and loop-unrolling constraints.
\item Explain invariants, asymptotic costs, and semantic failure modes for the studied transformations.
\end{itemize}

\section{Guided study sequence}
\begin{enumerate}
\item Read the named sections (8.2--8.7; pp. 379--444), then make a one-page concept map of local, regional, global and interprocedural optimization.
\item Write the input/output contract for every artifact introduced in the chapter.
\item Reconstruct a small original example by hand, recording the intermediate states.
\item Implement or inspect one corresponding lab and state two counterexamples.
\item Review a peer implementation: separate semantic correctness, compile-time cost, and runtime cost.
\end{enumerate}

\section{Worked design prompt (original)}
\subsection{Problem statement}
Simplify t1=a+b; t2=a+b; return t2, then add a potentially aliasing store.
\subsection{Required method}
\begin{itemize}
\item State assumptions explicitly; define token, type, IR, or target semantics as applicable.
\item Carry out each derivation step and annotate the exact invariant it preserves.
\item Provide at least one negative case that a naive algorithm mishandles.
\item Finish with a test oracle independent from the implementation.
\end{itemize}


\section{Laboratory bridge}
Implement expression-local value numbering with side-effect barriers.
\evidence{Before/after IR, equivalence tests, performance protocol and invalidation list.}

\section{Oral examination and failure analysis}
\begin{enumerate}
\item Name an optimization that speeds hot code while making total compilation more expensive.
\item What input changes would invalidate an existing cached result?
\item How should a diagnostic identify both the violated rule and the relevant source location?
\item What must be measured before claiming an implementation is faster?
\end{enumerate}

\section{Mastery rubric}
\begin{longtable}{|p{0.17\linewidth}|p{0.36\linewidth}|p{0.40\linewidth}|}
\hline
Dimension & Minimum evidence & Advanced evidence\\ \hline\endfirsthead
\hline
Dimension & Minimum evidence & Advanced evidence\\ \hline\endhead
Semantics & A terminating, correct reference test & Differential and metamorphic tests\\ \hline
Formal model & Definitions and one valid proof outline & Counterexample to an invalid assumption\\ \hline
Engineering & Reproducible build and trace & Instrumented cost/performance and review\\ \hline
\end{longtable}

\section{Source versus contemporary practice}
\textbf{Book-derived scope:} The chapter number, heading and topic references above follow the 2023 third edition.\par\modern Security limits, reproducibility protocols, differential fuzzing, SBOMs, release gates, and current tool integrations elsewhere in this suite are independent additions; they are not asserted to be practices advocated in the book.

\end{document}
